% Figure: heat loss per metre from the insulated steam pipe versus
% insulation thickness, the diminishing-returns curve from the worked
% example. Data points at 0, 25, 50, 100 mm from the chapter table, with a
% smooth interpolating curve. Annotates the diminishing marginal benefit.
\begin{figure}[htbp]
\centering
\begin{tikzpicture}
\begin{axis}[
  width=12cm, height=7.5cm,
  xlabel={insulation thickness $t$ (mm)},
  ylabel={heat loss per metre $\dot{Q}/L$ (W/m)},
  xmin=0, xmax=110, ymin=0, ymax=210,
  axis lines=left,
  legend pos=north east,
  legend cell align=left,
  tick label style={font=\scriptsize},
  label style={font=\small},
  legend style={font=\scriptsize},
  clip=false,
]
% chapter table values
\addplot[engblue, very thick, mark=*, mark options={fill=engblue}]
  coordinates {(0,197) (25,71) (50,49) (100,32)};
\addlegendentry{$\dot{Q}/L$ vs thickness}
% annotate the steep first increment and the flat tail
\draw[engred, ->, thick] (axis cs:6,150) -- (axis cs:6,90);
\node[font=\scriptsize, engred, anchor=west] at (axis cs:8,160)
  {first $25$ mm removes $64\%$ of the loss};
\draw[enggray, dashed] (axis cs:50,0) -- (axis cs:50,49);
\node[font=\scriptsize, enggray, anchor=south west] at (axis cs:52,52)
  {diminishing returns};
\end{axis}
\end{tikzpicture}
\caption[Heat loss versus insulation thickness for the steam pipe]{Heat
loss per metre from the insulated steam pipe of the worked example, as a
function of insulation thickness. The first $25\,$mm of insulation
removes nearly two thirds of the bare-pipe loss; each further increment
removes less, because the added conductive resistance is a shrinking
fraction of a total that already includes the earlier layers and the
outer-film resistance. The economic thickness sits where the marginal
value of the heat saved over the system lifetime equals the marginal cost
of insulation, well to the left of the point of vanishing returns. For
this pipe the outer radius is far above the critical insulation radius
$r_c=k/h$, so insulation always reduces the loss; on a thin wire the
opposite can hold.}
\label{fig:vol04ch05:insulation-thickness}
\end{figure}
